\begin{abstract}
Direct native-resolution comparison with 300 m FCOVER can confound retrieval disagreement with target-grid mismatch. We harmonized Sentinel-2, Landsat 8/9, and MODIS NDVI to the native FCOVER grid, replicated the protocol across four deliberately selected contrasting Qinghai Plateau areas of interest (AOIs), and combined block-grouped diagnostics with chronological historical-window evaluation. The 2021--2025 design comprised 144 formal OLS runs and 48 target-year NDVI-based dimidiate pixel model (DPM) configurations. Mean 2025 Multi-AOI OLS RMSE was 0.041023 for Sentinel-2, 0.036840 for Landsat, and 0.030787 for MODIS (unweighted means of 24 runs per sensor). Agreement, preferred histories, and coefficients varied strongly by AOI; these summaries do not provide cross-sensor inferential rankings. Selected OLS errors were lower than selected DPM errors in all 12 sensor--AOI comparisons, but this secondary benchmark used empirical target-year NDVI endpoints and retrospectively selected minima. All 72 Rolling-Origin runs completed; trajectories and paired 5 km block contrasts showed no uniform association between longer historical-window configurations and lower forward error. Primary geographic and historical-transfer interpretations were qualitatively robust to aggregation-order, non-overlapping temporal-composition, and Landsat aerosol-QA sensitivities. The common target grid removes a direct native-pixel-versus-FCOVER-grid comparison but not support differences after QA. Results quantify agreement with an operational FCOVER reference, not field-level FVC accuracy.
\end{abstract}
\noindent\textbf{Keywords:} fractional vegetation cover; NDVI; Copernicus FCOVER; common target grid; Rolling-Origin evaluation; spatial validation
